% ════════════════════════════════════════════════════════════════
\section{Wykład 6}
\subsection{Algorytmy Zachłanne}
% ════════════════════════════════════════════════════════════════

\textbf{Pomysł:} Algorytm w każdym kroku dokonuje lokalnie najlepszego wyboru (najlepszego w danym momencie).

\paragraph{Problem wyboru zajęć}
$S = \{z_1, \dots, z_n\}$ -- zbiór zajęć wymagający tych samych zasobów.
\begin{flalign*}
	s_i & \text{ -- czas rozpoczęcia zajęcia } z_i &  & s_i, t_i \in \N, \quad s_i < t_i &  & \\
	t_i & \text{ -- czas zakończenia zajęcia } z_i &  &                                  &  &
\end{flalign*}

Zajęcia $z_i$ i $z_j$ nazywamy \underline{zgodnymi}, jeśli $\langle s_i, t_i) \cap \langle s_j, t_j) = \emptyset$.
Zadanie polega na znalezieniu podzbioru o największej liczności $A \subseteq S$ parami zgodnych zajęć.

\begin{table}[H]
	\centering
	\renewcommand{\arraystretch}{1.2}
	\begin{NiceTabular}{c|cccccc}[margin,hvlines]
		$i$   & 1 & 2 & 3 & 4  & 5 & 6 \\
		\hline
		$s_i$ & 0 & 2 & 4 & 7  & 1 & 7 \\
		$t_i$ & 4 & 6 & 7 & 10 & 3 & 9
	\end{NiceTabular}
\end{table}

Porządkujemy zajęcia w czasie niemalejącym względem czasu zakończenia. To zajmuje $\BO(n \log n)$ czasu. Od tej pory zakładamy $t_1 \leq t_2 \leq \dots \leq t_n$ (zajęcia są już posortowane).

\begin{table}[H]
	\centering
	\renewcommand{\arraystretch}{1.2}
	\begin{NiceTabular}{c|cccccc}[margin,hvlines]
		$i$   & 1 & 2 & 3 & 4 & 5 & 6  \\
		\hline
		$s_i$ & 1 & 0 & 2 & 4 & 7 & 7  \\
		$t_i$ & 3 & 4 & 6 & 7 & 9 & 10
	\end{NiceTabular}
\end{table}

\paragraph{Podejście programowania dynamicznego}
Definiujemy zbiory zajęć:
\begin{flalign*}
	S_{ij} = \{ z_k \in S : t_i \leq s_k < t_k \leq s_j \} &  &  &  &
\end{flalign*}
(zbiór zajęć zaczynających się po czasie zakończenia zajęcia $z_i$ i kończących się przed czasem rozpoczęcia zajęcia $z_j$).

Dodajemy 2 sztuczne zajęcia: $z_0$ z czasem $t_0 = 0$ i $z_{n+1}$ z czasem $s_{n+1} = \infty$. Wtedy $S = S_{0, n+1}$.
Jeśli $i \geq j$, to $S_{ij} = \emptyset$, bo $s_j < t_j \leq t_i$ (bo zajęcia są posortowane).

Problem: znaleźć podzbiór o największej liczności zbioru $S_{ij}$ składający się z parami zgodnych zajęć dla $0 \leq i \leq j \leq n+1$.

Niech $A_{ij}$ będzie jakimkolwiek podzbiorem największej liczności składającym się z parami zgodnych zajęć zbioru $S_{ij}$ oraz niech $c(i, j) = |A_{ij}|$.

\paragraph{Podproblem $S_{ij}$:}
Załóżmy, że $z_k \in S_{ij}$. Jeśli wybierzemy $z_k$ do zbioru rozwiązania, to otrzymamy dwa zbiory (podproblemy):
\begin{itemize}
	\item $S_{ik}$ (zbiór zajęć, które zaczynają się po zakończeniu zajęcia $z_i$ i kończą się przed rozpoczęciem zajęcia $z_k$),
	\item $S_{kj}$ (zbiór zajęć, które zaczynają się po zakończeniu zajęcia $z_k$ i kończą przed rozpoczęciem zajęcia $z_j$).
\end{itemize}
Rozwiązanie dla $S_{ij}$ będzie sumą rozwiązań dla $S_{ik}$ oraz $S_{kj}$ oraz elementu $\{z_k\}$.

\paragraph{Własność optymalnej podstruktury:}
Załóżmy, że $z_k \in A_{ij}$. Wtedy $A_{ij} = A_{ik} \cup \{z_k\} \cup A_{kj}$. Gdyby istniało rozwiązanie $A'_{ik}$ dla $S_{ik}$ takie, że $|A'_{ik}| > |A_{ik}|$, to zbiór $A'_{ij} = A'_{ik} \cup \{z_k\} \cup A_{kj}$ byłby rozwiązaniem dla $S_{ij}$ takim, że $|A'_{ij}| > |A_{ij}|$. Sprzeczność. Analogicznie dla $S_{kj}$.
Zatem $A_{ik}$ oraz $A_{kj}$ muszą być rozwiązaniami optymalnymi dla odpowiednio $S_{ik}$ i $S_{kj}$.

Stąd wzór rekurencyjny:
\begin{flalign*}
	 & \begin{cases}
		   c(i, j) = 0 & \text{dla } i \geq j \\
		   c(i, j) = \max_{i < k < j} \{ c(i, k) + c(k, j) + 1 \}
	   \end{cases} &  &  &  &
\end{flalign*}

Używając tego wzoru rekurencyjnego możemy skonstruować algorytm programowania dynamicznego. Mamy $\BO(n^2)$ podproblemów. Dla $j > i$ mamy $j - i - 1$ wartości $k$ do sprawdzenia, czyli rozwiązanie jednego podproblemu zajmie czas liniowy. Otrzymujemy algorytm o złożoności $\BO(n^3)$.

\paragraph{Algorytm zachłanny}
\begin{lemma}
	Niech $z_m \in S_{ij}$ będzie zajęciem takim, że $t_m = \min \{t_k : z_k \in S_{ij}\}$. Wtedy:
	\begin{enumerate}
		\item Istnieje podzbiór o największej liczności składający się z parami zgodnych zajęć zbioru $S_{ij}$ zawierający $z_m$.
		\item $S_{im} = \emptyset$ (zatem po wybraniu $z_m$ do zbioru rozwiązania zostaje co najwyżej 1 niepusty podproblem do rozwiązania -- $S_{mj}$).
	\end{enumerate}
\end{lemma}

\begin{proof}~\\
	\textbf{Ad 2.} Przypuśćmy, że nie. Niech $z_k \in S_{im}$. Wtedy $t_i \leq s_k < t_k \leq s_m < t_m \leq s_j$. Zatem $z_k \in S_{ij}$ oraz $t_k < t_m$ -- sprzeczność z definicją zajęcia $z_m$.

	\textbf{Ad 1.} Niech $A_{ij}$ będzie podzbiorem o największej liczności składającym się z zajęć parami zgodnych zbioru $S_{ij}$. Niech $z_k$ będzie zajęciem z $A_{ij}$ o najmniejszym czasie zakończenia.
	Jeśli $z_k = z_m$, to teza jest udowodniona.
	Przypuśćmy, że $z_k \neq z_m$ i rozważmy zbiór $A'_{ij} = A_{ij} \setminus \{z_k\} \cup \{z_m\}$. Skoro $t_m \leq t_k$ oraz zajęcia z $A_{ij}$ są parami zgodne, to zajęcia z $A'_{ij}$ również są parami zgodne.
	Mamy $|A_{ij}| = |A'_{ij}|$, więc $A'_{ij}$ jest rozwiązaniem optymalnym zawierającym $z_m$.
\end{proof}

Z lematu wynika, że nie musimy sprawdzać $j - i - 1$ wartości $k$, możemy wybrać zajęcie o najmniejszym czasie zakończenia $z_m \in S_{ij}$.
Żeby rozwiązać podproblem $S_{ij}$ możemy dodać $z_m$ do zbioru rozwiązania a następnie dodać do zbioru rozwiązania rozwiązanie optymalne podproblemu $S_{mj}$. Zaczynamy z $S = S_{0, n+1}$ i $t_0 = 0$. W pierwszym kroku algorytm wybiera $z_1$.

\begin{alg}{Zachłanny wybór zajęć}{alg:greedyactivity}
	\FunctionNN{GreedyActivitySelector}{$s, t$} \Comment{zajęcia są posortowane po czasie zakończeń}
	\State $n \Assign \Length{s}$
	\State $A \Assign \{z_1\}$ \Comment{$A$ - zbiór rozwiązań}
	\State $j \Assign 1$
	\For{$i \Assign 2$ \To $n$}
	\If{$s_i \geq t_j$}
	\State $A \Assign A \cup \{z_i\}$
	\State $j \Assign i$
	\EndIf
	\EndFor
	\State \Return $A$
	\EndFunction
\end{alg}

Algorytm zachłanny w każdym kroku wybiera zajęcie zgodne ze wszystkimi zajęciami poprzednio wybranymi, o najmniejszym czasie zakończenia, aby zostawić jak najwięcej miejsca dla zajęć następnych (algorytm zachłanny). Skoro one są posortowane, jeśli $j < k$ i $z_k$ jest zgodne z $z_j$, to $z_k$ jest też zgodne z $z_i$ dla każdego $1 \leq i < j$.

\paragraph{Analiza poprawności:}
Algorytm zwraca zbiór parami zgodnych zajęć, bo dodaje do zbioru rozwiązania zajęcie tylko jeśli jest zgodne ze wszystkimi zajęciami dodanymi wcześniej.

Pokażemy przez indukcję po $n$, że algorytm znajduje rozwiązanie optymalne.
\begin{itemize}
	\item Dla $n = 1$ oczywiste $\checkmark$
	\item Niech $n > 1$. Z lematu istnieje optymalne rozwiązanie $A_1$ takie, że $z_1 \in A_1$. \\
	      Rozważmy problem dla $S' = \{z_i \in S : s_i \geq t_1\}$. Oczywiście zachodzi $|S'| < |S|$. \\
	      Dla każdego $z_j \in A_1 \setminus \{z_1\}$, mamy $z_j \in S'$, bo $s_j \geq t_1$, co wynika z faktu, że $A_1$ jest rozwiązaniem dla $S$. \\
	      Pokażemy, że $A_1 \setminus \{z_1\}$ jest rozwiązaniem optymalnym dla $S'$. \\
	      Przypuśćmy, że istnieje rozwiązanie $B$ dla $S'$ takie, że $|B| > |A_1 \setminus \{z_1\}|$. Wtedy $\{z_1\} \cup B$ byłoby rozwiązaniem dla $S$ takim, że:
	      \begin{flalign*}
		      |\{z_1\} \cup B| = 1 + |B| > 1 + |A_1 \setminus \{z_1\}| = |A_1| &  &  &  &
	      \end{flalign*}
	      Otrzymujemy sprzeczność, bo $A_1$ jest z założenia rozwiązaniem optymalnym. \\
	      Z założenia indukcyjnego nasz algorytm znajduje rozwiązanie optymalne $B'$ dla $S'$. Wtedy $|B'| = |A_1 \setminus \{z_1\}| = |A_1| - 1$. \\
	      Zbiór $\{z_1\} \cup B'$ jest rozwiązaniem znalezionym przez algorytm dla $S$ i:
	      \begin{flalign*}
		      |\{z_1\} \cup B'| = 1 + |B'| = 1 + |A_1| - 1 = |A_1| &  &  &  &
	      \end{flalign*}
	      Więc $\{z_1\} \cup B'$ jest rozwiązaniem optymalnym.
\end{itemize}

\paragraph{Złożoność czasowa}
\begin{itemize}[nosep]
	\item linia 1 -- $\BO(n)$
	\item linia 2 -- $\BO(1)$
	\item iteracji pętli w linii 3 jest $n - 1$
	\item wnętrze pętli zajmuje czas $\BO(1)$
\end{itemize}

Całość zajmuje $\BT(n)$, jeśli zajęcia są już posortowane po czasie zakończenia.
Ostatecznie złożoność algorytmu wynosi (uwzględniając początkowe sortowanie):
\begin{flalign*}
	\BO(n \log n) + \BT(n) = \BO(n \log n) &  &  &  &
\end{flalign*}
